\subsection{From sampled values to a formula}
\label{subsec:rotated-regular-polygons-empirical-curves}

The formula was first suggested by a finite sweep of rotation angles.  At each
sampled angle we evaluated the Haim--Kislev quadratic programs from
Section~\ref{sec:quadratic-program-algorithm-hk2019} in floating-point
arithmetic.  The samples trace a symmetric curve with its maximum at the
midpoint angle \(\pi/10\); see Figure~\ref{fig:rotated-pentagons-curve}.  This
was useful for finding the candidate branch and formula, but it cannot exclude
an unsampled crossing with another branch.  The theorem below instead uses an
exhaustive finite reduction and exact algebraic comparisons.

\begin{figure}[h!]
  \centering
  \includegraphics[width=0.92\linewidth]{working/rotated-regular-polygons/lagrangian-products-5x5.png}
  \caption{Floating-point finite-QP samples for \(K_\theta\), shown as blue
  markers, together with the exact profile proved below.  The horizontal line
  marks the conjectured Viterbo bound.  The agreement of the samples with the
  exact curve is an illustration and discovery check, not an input to the
  proof.}
  \label{fig:rotated-pentagons-curve}
\end{figure}
